%% kapa: exam=1 round=37 year=2026 period=1 form=A subject=civil_law subject_ko=민법 seq=21 number=21
%% source: Q-Net 감정평가사 기출문제 (cst00309) artlSeq=5258924 file=raw/1008/5258924/2026년 제37회 감정평가사 1차 1교시.pdf
%% conversion: kordoc 4.24.51 -> tools/gen1.py (auto); review=proofread-sonnet; flags=none
\begin{question}{21}
물권의 객체에 관한 설명으로 옳은 것은? (다툼이 있으면 판례에 따름)
\begin{choices}
  \item 건물의 개수는 공부상 등록의 개수에 의한다.
  \item 특정 양식장 내의 뱀장어일지라도 성장을 계속하는 한, 그 뱀장어 전부에 대한 양도 담보계약은 그 담보목적물이 특정되지 않아 성립하지 않는다.
  \item 권원 없이 타인의 토지에 심은 수목은 독립한 물권의 객체가 될 수 없다.
  \item 물권의 객체는 물건에 한정된다.
  \item 토지의 소유권은 지표에 관하여 그 효력을 미칠뿐 토지의 상하에 미치지 않는다.
\end{choices}
\end{question}
\answer{3}
